\chapter{Discusión}
\label{chap:discusion}

\section{Qué aporta la estimación semiparamétrica}

La evidencia predictiva favorece a E3 frente a las referencias principales. La ventaja aparece en las cinco ventanas del \textit{walk-forward} y se mantiene en el test y en el periodo posterior. Sin embargo, la comparación con referencias reforzadas acota su magnitud: en el test, el mejor Hawkes recupera cerca del 74\,\% de la mejora sobre AS y el modelo lineal con las mismas covariables, alrededor del 90\,\%. El estado del mercado aporta, por tanto, la mayor parte de la información que falta a la intensidad constante.

La ventaja restante corresponde a la especificación flexible de E3, incluida su modulación de la curva. No puede atribuirse íntegramente a la no linealidad de $g$ mediante la comparación con el GLM, porque las especificaciones difieren también en otros componentes. Las ablaciones muestran además que las variables de barrido aportan poco a la puntuación global, pero ayudan a corregir el nivel de las predicciones tras esos eventos. Frente a E3-fkX, esa aportación global no es significativa en el ask del test ni de la simulación. La calibración sigue siendo una limitación: el criterio tras barridos compradores falla en el periodo posterior.

Estas conclusiones responden a P1 y a la estabilidad temporal de P4. Se refieren al número esperado de eventos que alcanzan una distancia. Un AUC alto en esa tarea no mide acierto en la dirección del precio, probabilidad de beneficio ni ejecución de una orden propia.

\section{De la intensidad a la política y al resultado económico}

La intensidad estimada sí modifica la política, lo que responde a P3. La curva de distancia y la asimetría entre lados cambian las cotizaciones HJB con coeficientes congelados. En el test, con cruce estricto, E3 tesis ejecuta entre un 31 y un 38\,\% menos que AS y Hawkes con la misma forma de curva, y pierde menos por día. Esta reducción no equivale a una ventaja por ejecución: las diferencias pareadas de PnL por nocional no son significativas en ese periodo.

La respuesta a P2 depende de la regla de ejecución. Bajo cruce estricto y FIFO aproximada, el PnL bruto agregado es negativo. Bajo toque optimista, la HJB con E3 obtiene ganancias brutas pequeñas y supera a AS y Hawkes; E3 base, que no es la mejor especificación predictiva, presenta el mejor resultado económico entre las variantes de E3. Ninguno cubre las comisiones evaluadas. Por tanto, no se establece una correspondencia entre el orden predictivo de las intensidades y el orden económico de las políticas.

\section{Cómo interpretar las pérdidas}

La pérdida bajo cruce estricto no se explica solo por las comisiones: aparece antes de aplicarlas. La descomposición de la Sección~\ref{sec:p2-origen} muestra que la captura respecto del mid de decisión es positiva, pero el movimiento asociado a las ejecuciones es mayor y de signo contrario. Cotizar más lejos aumenta el margen inicial y reduce las ejecuciones; no garantiza mejor calidad en las que sí ocurren. Esta regularidad es compatible con selección adversa bajo la regla evaluada.

La interpretación concuerda con \citet{GlostenMilgrom1985}, para quienes el spread compensa el riesgo de negociar con agentes mejor informados, y con la evidencia empírica reciente en microestructura de alta frecuencia: la disyuntiva entre probabilidad de ejecución y retorno posterior documentada por \citet{AlbersEtAl2025} en futuros de Bitcoin en Binance, la deriva condicional negativa tras la ejecución de una orden límite estudiada por \citet{DeLise2024}, y el deterioro de la estrategia de control estocástico que \citet{LalorSwishchuk2024} examinan al introducir ejecuciones adversas y probabilidades de fill no adverso más realistas. La precisión estadística y el valor económico responden, por tanto, a criterios distintos \citep{LeitchTanner1991}.

Conviene separar selección adversa y gestión del inventario \citep{Menkveld2013}. El markout describe el valor de cada ejecución a un horizonte fijo; el PnL diario incorpora toda la trayectoria de las posiciones. Su semejanza de signo no identifica qué parte exacta del PnL causa cada mecanismo. La evidencia presentada respalda un deterioro alrededor de las ejecuciones y una menor pérdida total al operar menos; no demuestra que ajustar únicamente la aversión al inventario elimine ese deterioro.

La persistencia de resultados similares en los dos periodos refuerza la descripción bajo cruce estricto, pero no constituye una ley del mercado. Los eventos están agrupados, los modelos comparten reglas y parte de los datos había sido consultada. La excepción de R1 es decisiva: el conjunto de ejecuciones admitidas puede cambiar el signo bruto. Por ello, no se concluye que toda provisión de liquidez en BTC/USDT sea necesariamente deficitaria ni que la prioridad real sea irrelevante.

La comparación con motores alternativos sirve como comprobación auxiliar en configuraciones de entrenamiento. No valida de forma conjunta las políticas finales, la latencia real, las cancelaciones ni la posición en la cola de una orden en producción. La contribución económica es delimitar lo que no se monetizó bajo estos supuestos, sin convertir una explicación descriptiva en una afirmación causal universal.

\section{Limitaciones}
\label{sec:limitaciones}

\paragraph{Independencia de la evaluación.}
El test del 4 al 20 de septiembre no es completamente ciego: la especificación original de E3 ya se había evaluado allí, y una política de pausa se ejecutó antes de una orden de espera, aunque sus resultados no se consultaron. Las especificaciones finales se congelaron antes de su evaluación, pero esto no recupera la independencia del periodo. El \textit{walk-forward} se realizó dentro del bloque de desarrollo; sus ajustes respetan el orden temporal, sin constituir una selección anidada e independiente de toda la especificación.

La primera parte del periodo posterior, del 21 de septiembre al 6 de octubre, se usó en evaluaciones predictivas y económicas de versiones anteriores. Sus resultados no se usaron para seleccionar ni ajustar las especificaciones finales, pero ese tramo no es ciego en sentido estricto. Del 7 al 26 de octubre solo se construyó y verificó la capa de datos; ningún modelo se evaluó allí antes de la simulación.

\paragraph{Estimación y versiones del ajuste.}
Los Hawkes se estimaron con un tercio de los días por costo computacional, aunque sus capas de nivel, forma y señales usan todos. E3 utiliza tres ciclos de \textit{backfitting}, sin demostrar convergencia conjunta. En la versión original no se comprobaba la convergencia del último ajuste de la curva y el gradiente de $m(X)$ omitía el recorte de $b^\top z$ a $\pm5$, activado en unos 30 de 5,8 millones de segundos. La corrección verificó la convergencia de ese paso con $g$ y $h$ fijos. La curva cambia principalmente a 15--20~bps y el test conserva, a tres decimales, las ventajas de log-verosimilitud y el AUC. El \textit{walk-forward} conserva el ajuste original. Esta comprobación no cuantifica todo el efecto de una eventual convergencia conjunta incompleta.

\paragraph{Alcance, ejecución y prioridad.}
La respuesta estadística mide alcance, no fills propios. Mejorar su predicción no implica mejorar la predicción de las ejecuciones que obtendría un agente. R2 usa profundidad visible o imputada, no cuenta cancelaciones delante de la orden y pierde la prioridad al recotizar; R1 supone ejecución completa al toque sin cola delante. Son escenarios de sensibilidad, no cotas demostradas del PnL real. La resolución de 100~ms y los veinte niveles no permiten reconstruir la prioridad efectiva. La sensibilidad de la HJB con E3 impide generalizar su resultado bruto a cualquier regla de cola.

\paragraph{Control y condiciones de operación.}
La HJB usa coeficientes congelados, una curva interpolada en nueve nodos y tablas numéricas; no incorpora la dinámica de $X_t$ como estado. GLFT añade una aproximación exponencial local. La evaluación final supone activación al inicio del segundo sin latencia de envío, no los aproximadamente 130~ms de red observados en la captura. El inventario terminal se marca al mid, sin liquidación ejecutada, impacto propio ni costo de salida. Las sensibilidades anteriores de motores y latencia no cubren todas las políticas finales ni carreras de cancelación en microsegundos \citep{BudishCramtonShim2015,AquilinaBudishONeill2022}.

\paragraph{Incertidumbre y generalización.}
Los conteos están anidados por distancia y la puntuación Poisson es compuesta; evalúa la media, no la distribución conjunta. El remuestreo por días recoge la dependencia intradía, pero no garantiza independencia entre días. La calibración tras barridos sigue siendo inestable. El estudio se refiere a un mercado y a unas comisiones concretas: no demuestra ausencia de valor en otros activos, tarifas, plataformas o mecanismos de ejecución.
